% =============================================================================
% EC8070 – RESEARCH PROJECT III: THESIS  (10 pages including front page + refs)
% AI-Based Optimized Energy Utilization System Using Edge Controllers
% =============================================================================
\documentclass[10pt,a4paper,twocolumn]{article}

% ---- Packages ---------------------------------------------------------------
\usepackage[utf8]{inputenc}
\usepackage[T1]{fontenc}
\usepackage{lmodern}
\usepackage[top=0.75in,bottom=0.75in,left=0.6in,right=0.6in]{geometry}
\usepackage{graphicx}
\usepackage{amsmath,amssymb}
\usepackage{booktabs}
\usepackage{array}
\usepackage{multirow}
\usepackage{tabularx}
\usepackage{float}
\usepackage{caption}
\usepackage{subcaption}
\usepackage{hyperref}
\usepackage{enumitem}
\usepackage{setspace}
\usepackage{titlesec}
\usepackage{xcolor}
\usepackage{tikz}
\usetikzlibrary{shapes.geometric, arrows.meta, positioning, fit, backgrounds, calc, patterns, decorations.markings}
\usepackage{pgfplots}
\pgfplotsset{compat=1.18}
\usepackage{pgf-pie}
\usepackage{fancyhdr}
\usepackage{dblfloatfix}
\usepackage{stfloats}

% ---- Compact spacing --------------------------------------------------------
\setlength{\parskip}{2pt plus 1pt minus 1pt}
\setlength{\parindent}{0.8em}
\titlespacing*{\section}{0pt}{8pt}{4pt}
\titlespacing*{\subsection}{0pt}{6pt}{3pt}
\titlespacing*{\subsubsection}{0pt}{4pt}{2pt}
\titleformat{\section}{\normalsize\bfseries}{\thesection.}{0.5em}{}
\titleformat{\subsection}{\small\bfseries}{\thesubsection}{0.5em}{}
\titleformat{\subsubsection}{\small\itshape}{\thesubsubsection}{0.5em}{}
\setlist{nosep, leftmargin=1.2em}
\captionsetup{font=small, labelfont=bf, skip=4pt}

% ---- Hyperref ----------------------------------------------------------------
\hypersetup{colorlinks=true, linkcolor=blue!70!black, citecolor=blue!70!black, urlcolor=blue!70!black}

% ---- Page style --------------------------------------------------------------
\pagestyle{fancy}
\fancyhf{}
\fancyhead[L]{\footnotesize AI-Based Optimized Energy Utilization System}
\fancyhead[R]{\footnotesize EC8070}
\fancyfoot[C]{\footnotesize\thepage}
\renewcommand{\headrulewidth}{0.4pt}

% ==============================================================================
% FRONT PAGE (Page 1)
% ==============================================================================
\begin{document}
\onecolumn
\begin{titlepage}
    \centering
    \vspace*{2cm}

    {\LARGE\bfseries AI-BASED OPTIMIZED ENERGY UTILIZATION SYSTEM \\[0.3cm]
    USING EDGE CONTROLLERS}

    \vspace{2cm}

    {\large THESIS}\\[0.4cm]
    {\normalsize IN PARTIAL FULFILMENT OF THE REQUIREMENTS OF THE DEGREE OF}\\[0.2cm]
    {\normalsize BACHELOR OF SCIENCE IN ENGINEERING}

    \vspace{2cm}

    {\normalsize Submitted By:}\\[0.5cm]
    % TODO: Replace with actual group member names and registration numbers
    {\normalsize MEMBER~ONE~A.B.C (2021/E/XXX)}\\[0.2cm]
    {\normalsize MEMBER~TWO~D.E.F (2021/E/YYY)}\\[0.2cm]
    {\normalsize MEMBER~THREE~G.H.I (2021/E/ZZZ)}

    \vfill

    {\normalsize DEPARTMENT OF COMPUTER ENGINEERING}\\[0.2cm]
    {\normalsize FACULTY OF ENGINEERING}\\[0.2cm]
    {\normalsize UNIVERSITY OF JAFFNA}\\[0.5cm]
    {\normalsize JULY 2026}
\end{titlepage}

% ---- Table of Contents (Page 2) ---------------------------------------------
\tableofcontents
\thispagestyle{fancy}

% ---- Switch to two-column for the rest (Pages 3–10) -------------------------
\newpage
\twocolumn

% ==============================================================================
% ABSTRACT
% ==============================================================================
\section*{Abstract}
\addcontentsline{toc}{section}{Abstract}

Rising electricity costs under the LECO Time-of-Use (TOU) tariff in Sri Lanka demand intelligent Home Energy Management Systems (HEMS). This thesis presents an AI-based HEMS that integrates LSTM power demand forecasting, an MQTT edge controller for real-time TOU pricing and grid capacity signals, a local LLM (Ollama / LLaMA~3.2) for natural-language user preference parsing, and Mixed-Integer Linear Programming (MILP) for globally optimal appliance scheduling---all orchestrated by a LangGraph state-machine agent. The LSTM model, trained on 10,080 one-minute power readings, predicts 24-hour demand for five domestic appliances. The MILP solver computes cost-minimizing binary ON/OFF schedules across 24~hourly slots subject to capacity, runtime, comfort, and user-preference constraints. Results show a \textbf{37.0\% daily cost reduction} (LKR~2,692.80~$\rightarrow$~1,697.10), with an 82-policy validation suite achieving 100\% pass rate covering capacity compliance, LSTM accuracy (93.3\% average), cost optimality, and weather-based comfort. The system provides a practical, reproducible framework for intelligent residential energy management in the Sri Lankan context.

\noindent\textbf{Keywords:} HEMS; LSTM forecasting; MILP; LLM; edge controller; TOU tariff; MQTT; demand-side management.

% ==============================================================================
% 1. INTRODUCTION
% ==============================================================================
\section{Introduction}

The residential sector accounts for a substantial share of electricity consumption in Sri Lanka. The Lanka Electricity Company (LECO) introduced Time-of-Use (TOU) domestic tariffs with three pricing bands---off-peak (LKR~33/kWh), day (LKR~47/kWh), and peak (LKR~106/kWh)---creating a threefold price differential~\cite{leco2024}. However, most households lack the automation tools to exploit these differentials.

Home Energy Management Systems (HEMS) have emerged as a solution for automated demand-side management~\cite{zhou2016}. While significant progress has been made in high-resource settings~\cite{beaudin2015, pipattanasomporn2012}, few implementations target Sri Lankan constraints including the LECO TOU structure, limited smart-grid infrastructure, and the need for edge-based processing.

\subsection{Research Problem}
The central problem is the absence of an integrated, AI-driven appliance scheduling system for Sri Lankan households that simultaneously: (i)~predicts appliance-level power demand, (ii)~respects real-time grid capacity constraints from edge controllers, (iii)~interprets natural-language user preferences, and (iv)~computes provably cost-optimal schedules.

\subsection{Aim and Objectives}
The aim is to design, implement, and validate an AI-based HEMS using edge controllers for optimized energy utilization under TOU pricing. The objectives are:
\begin{enumerate}[label=(\roman*)]
    \item Develop an LSTM model that predicts 24-hour appliance-level power demand.
    \item Design an MQTT edge controller interface for real-time TOU pricing and grid capacity.
    \item Integrate a local LLM for parsing natural-language preferences into scheduling constraints.
    \item Formulate and solve the scheduling problem as a MILP minimizing cost subject to capacity, runtime, comfort, and preference constraints.
    \item Validate through a comprehensive 82-policy test suite.
\end{enumerate}

% ==============================================================================
% 2. LITERATURE REVIEW
% ==============================================================================
\section{Literature Review}

\textbf{Demand Forecasting.} LSTM networks~\cite{hochreiter1997} have demonstrated superior performance in energy demand forecasting. Kong~et~al.~\cite{kong2019} applied LSTMs to household load forecasting with significant accuracy gains. Kim and Cho~\cite{kim2019} used CNN-LSTM hybrids for multi-step prediction.

\textbf{Optimization-Based Scheduling.} MILP is the dominant approach for appliance scheduling under TOU pricing~\cite{garifi2018}. Pipattanasomporn~et~al.~\cite{pipattanasomporn2012} achieved 21\% peak demand reduction using MILP-based HEMS. Beaudin and Zareipour~\cite{beaudin2015} identified MILP as the most reliable method for cost-optimal solutions.

\textbf{Edge Computing \& IoT.} MQTT has become standard for IoT communication in smart homes due to low overhead~\cite{naik2017}. Edge controllers enable local processing of grid signals, reducing latency and enabling operation during connectivity disruptions~\cite{zhou2016}.

\textbf{LLMs for Smart Homes.} Large language models have been explored for interpreting user intent in home automation~\cite{brown2020, king2023}. However, integrating LLMs for preference parsing within optimization-based HEMS remains unexplored.

\textbf{Gap Identification.} No prior work integrates LSTM forecasting, MILP scheduling, MQTT edge controllers, and LLM preference parsing within a single HEMS pipeline tailored to Sri Lankan TOU tariffs, with weather-aware comfort scheduling---the gap this thesis fills.

% ==============================================================================
% 3. METHODOLOGY
% ==============================================================================
\section{Methodology}

\subsection{Research Design}
We adopt a quantitative, experimental design with five phases: (1)~sensor data ingestion, (2)~LSTM forecasting, (3)~TOU/capacity communication via edge controllers, (4)~AI agent-based scheduling, and (5)~mobile app visualization. All experiments use deterministic random seeds for reproducibility.

\subsection{System Architecture}
The proposed system follows a layered edge-computing architecture (Fig.~\ref{fig:architecture}). Five domestic appliances publish power readings via MQTT. The LSTM predictor forecasts 24-hour demand. A LangGraph agent integrates TOU pricing, weather data, and LLM-parsed preferences to compute MILP-optimal schedules, delivered to a Flutter mobile app via Firebase Firestore.

% ---- SYSTEM ARCHITECTURE DIAGRAM ----
\begin{figure}[H]
    \centering
    \resizebox{\columnwidth}{!}{%
    \begin{tikzpicture}[
        node distance=0.8cm and 1.0cm,
        block/.style={rectangle, draw, rounded corners=3pt, minimum width=2.2cm, minimum height=0.7cm, align=center, font=\scriptsize\sffamily},
        arrow/.style={-{Stealth[length=2mm]}, thick, color=gray!70!black},
    ]
        \node[block, fill=blue!15] (sensor) {Smart Home\\5 Appliances};
        \node[block, fill=purple!15, right=1.2cm of sensor] (mqtt) {MQTT Broker\\Mosquitto};
        \node[block, fill=orange!15, below=of mqtt] (lstm) {LSTM Predictor\\Run\_LSTM.py};
        \node[block, fill=yellow!15, below=of lstm] (agent) {LangGraph Agent\\agent.py};
        \node[block, fill=green!15, left=of agent] (edge) {Edge Controller\\TOU+Capacity};
        \node[block, fill=cyan!15, right=of agent] (weather) {Open-Meteo\\Weather API};
        \node[block, fill=red!10, below left=0.6cm and 0cm of agent] (llm) {Ollama LLM\\LLaMA 3.2};
        \node[block, fill=teal!15, below right=0.6cm and 0cm of agent] (firebase) {Firebase\\Firestore};
        \node[block, fill=pink!15, below=1.8cm of agent] (app) {Flutter\\Mobile App};

        \draw[arrow] (sensor) -- node[above, font=\tiny]{MQTT} (mqtt);
        \draw[arrow] (mqtt) -- node[right, font=\tiny]{home/power} (lstm);
        \draw[arrow] (lstm) -- node[right, font=\tiny]{JSON} (agent);
        \draw[arrow] (edge) -- node[above, font=\tiny]{MQTT} (agent);
        \draw[arrow] (weather) -- node[above, font=\tiny]{HTTP} (agent);
        \draw[arrow] (agent) -- (llm);
        \draw[arrow] (agent) -- (firebase);
        \draw[arrow] (firebase) -- (app);
    \end{tikzpicture}}
    \caption{System architecture of the proposed AI-based HEMS.}
    \label{fig:architecture}
\end{figure}

\subsection{LSTM Power Demand Forecasting}
An LSTM neural network predicts next-hour power consumption for each appliance. The model consists of two stacked LSTM layers (64 hidden units each) with a dense output layer. It is trained on 10,080 one-minute power readings from five appliances (WashingMachine, Heater, AC, VehicleCharger, VacuumCleaner) collected via MQTT.

During operation, a rolling buffer of 1,464 samples ($\approx$24~hours) is maintained. Every 30 new MQTT messages, the model predicts and accumulates results into a 1,440-sample daily store. Predictions are processed into 24 hourly averages and binarized using a demand-driven approach:
\begin{equation}
    R_a = \text{round}\!\left(\frac{\sum_{m=1}^{1440} \hat{P}_{a,m}}{P_a^{\text{rated}}}\right)
    \label{eq:binarize}
\end{equation}
where $\hat{P}_{a,m}$ is the predicted power at minute~$m$ and $P_a^{\text{rated}}$ is the rated power. The top $R_a$ hours with highest predicted power are marked ON. This threshold-free approach avoids arbitrary cutoffs. Fig.~\ref{fig:lstm_pipeline} shows the LSTM pipeline.

% ---- LSTM PIPELINE FLOWCHART ----
\begin{figure}[H]
    \centering
    \resizebox{\columnwidth}{!}{%
    \begin{tikzpicture}[
        node distance=0.55cm,
        block/.style={rectangle, draw, rounded corners=2pt, minimum width=2.8cm, minimum height=0.55cm, align=center, font=\scriptsize\sffamily, fill=#1},
        arrow/.style={-{Stealth[length=1.8mm]}, thick, gray!70!black},
    ]
        \node[block=blue!10] (buf) {Rolling Buffer\\(1464 samples)};
        \node[block=orange!10, below=of buf] (scale) {MinMax Scaling\\(scaler.pkl)};
        \node[block=yellow!10, below=of scale] (seq) {Sequence Windows\\(24-step)};
        \node[block=red!10, below=of seq] (lstm) {LSTM Model\\(2$\times$64 units)};
        \node[block=green!10, below=of lstm] (inv) {Inverse Transform\\$\rightarrow$ Watts};
        \node[block=cyan!10, below=of inv] (avg) {24h Hourly\\Averages};
        \node[block=purple!10, below=of avg] (bin) {Demand-Driven\\Binarization};
        \node[block=teal!10, below=of bin] (out) {appliance\_data.json\\(ON/OFF states)};

        \draw[arrow] (buf) -- (scale);
        \draw[arrow] (scale) -- (seq);
        \draw[arrow] (seq) -- (lstm);
        \draw[arrow] (lstm) -- (inv);
        \draw[arrow] (inv) -- (avg);
        \draw[arrow] (avg) -- (bin);
        \draw[arrow] (bin) -- (out);
    \end{tikzpicture}}
    \caption{LSTM prediction pipeline: from sensor buffer to binary ON/OFF states.}
    \label{fig:lstm_pipeline}
\end{figure}

\subsection{Edge Controller \& TOU Communication}
The edge controller (\texttt{publish\_tou\_test.py}) scrapes live LECO TOU tariff data from the official website and publishes it with grid capacity limits via MQTT on topic \texttt{power/tou\_domestic} every 5~minutes. Capacity limits are configurable via \texttt{capacity\_config.json} (Table~\ref{tab:tou}).

\begin{table}[H]
    \centering
    \caption{LECO TOU tariff bands and grid capacity limits.}
    \label{tab:tou}
    \small
    \begin{tabular}{@{}lccc@{}}
        \toprule
        \textbf{Band} & \textbf{Time} & \textbf{Rate} & \textbf{Capacity} \\
        \midrule
        Off-Peak & 22:00--05:00 & 33 LKR/kWh & 6.0 kW \\
        Day      & 05:00--18:00 & 47 LKR/kWh & 4.0 kW \\
        Peak     & 18:00--22:00 & 106 LKR/kWh & 2.5 kW \\
        \bottomrule
    \end{tabular}
\end{table}

\subsection{LLM-Based User Preference Parsing}
User preferences are submitted as natural language via a Flutter app, stored in Firebase Firestore, and sent to Ollama (LLaMA~3.2, temperature~0.0) with a structured prompt. The LLM returns JSON specifying \texttt{allow\_peak} (boolean per appliance) and \texttt{preferred\_hours} (hour indices). Example: \textit{``Heater need to turn only night times''} $\rightarrow$ \texttt{preferred\_hours[Heater]}~=~\{18--23\}, \texttt{allow\_peak[Heater]}~=~false. If the LLM is unavailable, the system defaults to no peak access.

\subsection{MILP Optimal Scheduling}
The scheduling problem is formulated as a MILP with 120 binary decision variables ($x_{a,t} \in \{0,1\}$ for 5~appliances~$\times$~24~hours) plus 24 slack variables for soft capacity constraints.

\textbf{Objective:} Minimize total cost with comfort bias:
\begin{equation}
    \min \sum_{a \in \mathcal{A}} \sum_{t=0}^{23} P_a \cdot c_{a,t} \cdot x_{a,t} + 10^6 \sum_{t=0}^{23} s_t
    \label{eq:objective}
\end{equation}
where $c_{a,t}$ incorporates TOU price and a $-100$ comfort bias for AC during hot hours ($\geq$28\textdegree{}C or $\geq$80\% humidity) and Heater during cold hours ($\leq$20\textdegree{}C).

\textbf{Constraints:}
\begin{align}
    \sum_{t=0}^{23} x_{a,t} &= R_a \quad \forall\, a \in \mathcal{A} \label{eq:runtime} \\
    \sum_{a \in \mathcal{A}} P_a \cdot x_{a,t} - s_t &\leq C_t \quad \forall\, t \label{eq:capacity} \\
    x_{a,t} &= 0 \;\; \text{if blocked by user prefs} \label{eq:userprefs}
\end{align}

The MILP is solved using \texttt{scipy.optimize.milp} with the HiGHS backend, yielding a globally optimal binary schedule.

Table~\ref{tab:appliances} lists the five managed appliances with their rated power and key scheduling constraints.

\begin{table}[H]
    \centering
    \caption{Managed appliances and constraints.}
    \label{tab:appliances}
    \small
    \begin{tabular}{@{}lcc@{}}
        \toprule
        \textbf{Appliance} & \textbf{Power} & \textbf{Key Constraint} \\
        \midrule
        WashingMachine & 0.6 kW & Flexible, off-peak \\
        Heater         & 2.0 kW & Cold hours (comfort) \\
        AC             & 1.2 kW & Hot/humid (comfort) \\
        VehicleCharger & 2.2 kW & Overnight preferred \\
        VacuumCleaner  & 1.1 kW & Cost minimization \\
        \bottomrule
    \end{tabular}
\end{table}

\subsection{LangGraph Agent Workflow}
The agent is a 4-node LangGraph StateGraph executing every 30~minutes (Fig.~\ref{fig:langgraph}):

% ---- LANGGRAPH FLOWCHART ----
\begin{figure}[H]
    \centering
    \resizebox{\columnwidth}{!}{%
    \begin{tikzpicture}[
        node distance=0.7cm,
        block/.style={rectangle, draw, rounded corners=3pt, minimum width=2.6cm, minimum height=0.65cm, align=center, font=\scriptsize\sffamily, fill=#1},
        decision/.style={diamond, draw, aspect=2, minimum width=1.2cm, minimum height=0.5cm, align=center, font=\tiny\sffamily, fill=yellow!20},
        arrow/.style={-{Stealth[length=1.8mm]}, thick, gray!70!black},
        darrow/.style={-{Stealth[length=1.8mm]}, thick, red!60!black, dashed},
    ]
        \node[block=blue!15] (start) {START};
        \node[block=green!15, below=of start] (fetch) {\textbf{1. Fetch Data}\\MQTT + Weather + Firebase};
        \node[decision, below=of fetch] (d1) {Error?};
        \node[block=purple!15, below=of d1] (parse) {\textbf{2. Parse Preferences}\\Ollama LLM};
        \node[decision, below=of parse] (d2) {Error?};
        \node[block=orange!15, below=of d2] (milp) {\textbf{3. MILP Scheduling}\\scipy.optimize.milp};
        \node[block=teal!15, below=of milp] (write) {\textbf{4. Write Results}\\JSON + Firestore};
        \node[block=gray!15, below=of write] (endnode) {END};
        \node[block=red!15, right=1.5cm of milp] (fallback) {\textbf{Fallback}\\Revert to LSTM};

        \draw[arrow] (start) -- (fetch);
        \draw[arrow] (fetch) -- (d1);
        \draw[arrow] (d1) -- node[left, font=\tiny]{No} (parse);
        \draw[darrow] (d1.east) -| node[above, font=\tiny]{Yes} (fallback);
        \draw[arrow] (parse) -- (d2);
        \draw[arrow] (d2) -- node[left, font=\tiny]{No} (milp);
        \draw[darrow] (d2.east) -| node[above, font=\tiny]{Yes} (fallback);
        \draw[arrow] (milp) -- (write);
        \draw[arrow] (write) -- (endnode);
        \draw[darrow] (fallback) |- (endnode);
    \end{tikzpicture}}
    \caption{LangGraph agent workflow with conditional error handling.}
    \label{fig:langgraph}
\end{figure}

\subsection{Validation Framework}
An 82-policy deterministic test suite validates system correctness across seven categories (Table~\ref{tab:policies}). Reproducibility is ensured by fixed seeds, deterministic MILP solving, centralized configuration (\texttt{config.py}), and publicly available model weights.

\begin{table}[H]
    \centering
    \caption{Validation policy categories.}
    \label{tab:policies}
    \small
    \begin{tabular}{@{}lc@{}}
        \toprule
        \textbf{Category} & \textbf{Policies} \\
        \midrule
        Capacity (load $\leq$ grid/hour) & 48 \\
        Format (binary, 24 elements) & 10 \\
        LLM preference enforcement & 10 \\
        LSTM accuracy ($\geq$50\% coverage) & 5 \\
        Cost optimality (opt $\leq$ baseline) & 6 \\
        Savings threshold ($\geq$20\%) & 1 \\
        Weather comfort (AC/Heater) & 2 \\
        \midrule
        \textbf{Total} & \textbf{82} \\
        \bottomrule
    \end{tabular}
\end{table}

% ==============================================================================
% 4. RESULTS
% ==============================================================================
\section{Results}

\subsection{LSTM Prediction Accuracy}
The LSTM achieves 93.3\% average scheduling accuracy (Table~\ref{tab:lstm}). The Heater scheduled 2~hours instead of 3 because the LLM restricted it to night hours \{18--23\}, and the MILP allocated the 2~cheapest off-peak slots within that window.

\begin{table}[H]
    \centering
    \caption{LSTM accuracy: required vs.\ scheduled hours.}
    \label{tab:lstm}
    \small
    \begin{tabular}{@{}lccc@{}}
        \toprule
        \textbf{Appliance} & \textbf{Req.} & \textbf{Sched.} & \textbf{Match} \\
        \midrule
        WashingMachine & 7  & 7  & 100\% \\
        Heater         & 3  & 2  & 66.7\% \\
        AC             & 20 & 20 & 100\% \\
        VehicleCharger & 4  & 4  & 100\% \\
        VacuumCleaner  & 3  & 3  & 100\% \\
        \midrule
        \textbf{Average} & & & \textbf{93.3\%} \\
        \bottomrule
    \end{tabular}
\end{table}

\subsection{Cost Optimization Results}
The MILP achieves a \textbf{37.0\% total cost reduction} (LKR~2,692.80~$\rightarrow$~1,697.10). Table~\ref{tab:cost} and Fig.~\ref{fig:cost_bar} detail per-appliance savings.

\begin{table}[H]
    \centering
    \caption{Cost savings: baseline vs.\ optimized schedule.}
    \label{tab:cost}
    \small
    \begin{tabular}{@{}lrrr@{}}
        \toprule
        \textbf{Appliance} & \textbf{Base} & \textbf{Opt.} & \textbf{Save\%} \\
        \midrule
        WashingMachine & 216.0  & 155.4  & 28.1\% \\
        Heater         & 372.0  & 132.0  & \textbf{64.5\%} \\
        AC             & 1273.2 & 1010.4 & 20.6\% \\
        VehicleCharger & 611.6  & 290.4  & \textbf{52.5\%} \\
        VacuumCleaner  & 220.0  & 108.9  & 50.5\% \\
        \midrule
        \textbf{Total} & \textbf{2692.8} & \textbf{1697.1} & \textbf{37.0\%} \\
        \bottomrule
    \end{tabular}
\end{table}

% ---- COST COMPARISON BAR CHART ----
\begin{figure}[H]
    \centering
    \begin{tikzpicture}
        \begin{axis}[
            ybar,
            bar width=5pt,
            width=\columnwidth,
            height=4.2cm,
            ylabel={\scriptsize Cost (LKR)},
            ylabel style={at={(axis description cs:0.08,0.5)}},
            symbolic x coords={WM, Heater, AC, VC, Vacuum},
            xtick=data,
            x tick label style={font=\tiny},
            y tick label style={font=\tiny},
            ymin=0, ymax=1400,
            legend style={at={(0.5,-0.25)}, anchor=north, legend columns=2, font=\tiny},
            nodes near coords,
            nodes near coords style={font=\tiny, above, rotate=45},
            every axis plot/.append style={fill opacity=0.85},
            enlarge x limits=0.15,
        ]
            \addplot[fill=red!60] coordinates {
                (WM, 216) (Heater, 372) (AC, 1273) (VC, 612) (Vacuum, 220)
            };
            \addplot[fill=green!60] coordinates {
                (WM, 155) (Heater, 132) (AC, 1010) (VC, 290) (Vacuum, 109)
            };
            \legend{Baseline, Optimized}
        \end{axis}
    \end{tikzpicture}
    \caption{Per-appliance cost comparison (LKR).}
    \label{fig:cost_bar}
\end{figure}

\subsection{Schedule Shift Analysis}
The optimizer shifts appliance hours from expensive to cheap TOU bands:
\begin{itemize}
    \item \textbf{WashingMachine:} Hours 13, 15, 18 (day/peak) $\rightarrow$ 00, 01, 03 (off-peak).
    \item \textbf{Heater:} Hours 04, 16, 21 $\rightarrow$ 22, 23 (off-peak, night window).
    \item \textbf{AC:} Peak hours 18--20 $\rightarrow$ 01, 22, 23 (off-peak).
    \item \textbf{VehicleCharger:} Hours 20--23 $\rightarrow$ 00--04 (overnight off-peak).
    \item \textbf{VacuumCleaner:} Hours 05, 16, 19 $\rightarrow$ 01, 22, 23 (off-peak).
\end{itemize}

Fig.~\ref{fig:heatmap} visualizes the original vs.\ optimized schedules as a 24-hour heatmap.

% ---- 24-HOUR SCHEDULE HEATMAP (Full width at bottom of page) ----
\begin{figure*}[b]
    \centering
    \begin{tikzpicture}[scale=0.42]
        % Data: Original states
        % WM:  0 0 1 0 1 0 1 1 0 0 0 0 0 1 0 1 0 0 1 0 0 0 0 0
        % HT:  0 0 0 0 1 0 0 0 0 0 0 0 0 0 0 0 1 0 0 0 0 1 0 0
        % AC:  1 0 1 1 1 1 1 1 1 1 1 1 1 1 1 1 1 1 1 1 1 0 0 0
        % VC:  0 0 0 0 0 0 0 0 0 0 0 0 0 0 0 0 0 0 0 0 1 1 1 1
        % VV:  0 0 0 0 0 1 0 0 0 0 0 0 0 0 0 0 1 0 0 1 0 0 0 0

        % Data: Optimized states
        % WM:  1 1 1 1 1 0 1 1 0 0 0 0 0 0 0 0 0 0 0 0 0 0 0 0
        % HT:  0 0 0 0 0 0 0 0 0 0 0 0 0 0 0 0 0 0 0 0 0 0 1 1
        % AC:  1 1 1 1 1 1 1 1 1 1 1 1 1 1 1 1 1 1 0 0 0 0 1 1
        % VC:  1 0 1 1 1 0 0 0 0 0 0 0 0 0 0 0 0 0 0 0 0 0 0 0
        % VV:  0 1 0 0 0 0 0 0 0 0 0 0 0 0 0 0 0 0 0 0 0 0 1 1

        % Title
        \node[font=\footnotesize\bfseries] at (12, 12.5) {Original Schedule (LSTM Predicted)};
        \node[font=\footnotesize\bfseries] at (38, 12.5) {Optimized Schedule (MILP)};

        % TOU band coloring helper - Original side
        % Off-peak: 0-4, 22-23  Day: 5-17  Peak: 18-21
        \foreach \h in {0,...,4} {
            \fill[green!8] (\h, -0.5) rectangle (\h+1, 5.5);
        }
        \foreach \h in {5,...,17} {
            \fill[yellow!8] (\h, -0.5) rectangle (\h+1, 5.5);
        }
        \foreach \h in {18,...,21} {
            \fill[red!8] (\h, -0.5) rectangle (\h+1, 5.5);
        }
        \foreach \h in {22,...,23} {
            \fill[green!8] (\h, -0.5) rectangle (\h+1, 5.5);
        }

        % TOU band coloring helper - Optimized side (offset by 26)
        \foreach \h in {0,...,4} {
            \pgfmathsetmacro{\x}{\h+26}
            \fill[green!8] (\x, -0.5) rectangle (\x+1, 5.5);
        }
        \foreach \h in {5,...,17} {
            \pgfmathsetmacro{\x}{\h+26}
            \fill[yellow!8] (\x, -0.5) rectangle (\x+1, 5.5);
        }
        \foreach \h in {18,...,21} {
            \pgfmathsetmacro{\x}{\h+26}
            \fill[red!8] (\x, -0.5) rectangle (\x+1, 5.5);
        }
        \foreach \h in {22,...,23} {
            \pgfmathsetmacro{\x}{\h+26}
            \fill[green!8] (\x, -0.5) rectangle (\x+1, 5.5);
        }

        % Y-axis labels
        \node[font=\tiny, anchor=east] at (-0.2, 4.5) {WM};
        \node[font=\tiny, anchor=east] at (-0.2, 3.5) {Heater};
        \node[font=\tiny, anchor=east] at (-0.2, 2.5) {AC};
        \node[font=\tiny, anchor=east] at (-0.2, 1.5) {VCharger};
        \node[font=\tiny, anchor=east] at (-0.2, 0.5) {Vacuum};

        % X-axis labels - Original
        \foreach \h in {0,2,...,22} {
            \node[font=\tiny] at (\h+0.5, -1) {\h};
        }
        % X-axis labels - Optimized
        \foreach \h in {0,2,...,22} {
            \pgfmathsetmacro{\x}{\h+26}
            \node[font=\tiny] at (\x+0.5, -1) {\h};
        }

        % X-axis title
        \node[font=\tiny] at (12, -1.8) {Hour of Day};
        \node[font=\tiny] at (38, -1.8) {Hour of Day};

        % Draw grid - Original (24 cols x 5 rows)
        \draw[very thin, gray] (0, -0.5) grid[step=1] (24, 5.5);

        % ORIGINAL SCHEDULE - fill ON cells
        % WM row (y=4): hours 2,4,6,7,13,15,18
        \foreach \h in {2,4,6,7,13,15,18} {
            \fill[blue!50] (\h+0.05, 4.05) rectangle (\h+0.95, 4.95);
        }
        % Heater row (y=3): hours 4,16,21
        \foreach \h in {4,16,21} {
            \fill[red!50] (\h+0.05, 3.05) rectangle (\h+0.95, 3.95);
        }
        % AC row (y=2): hours 0,2-20
        \foreach \h in {0,2,3,4,5,6,7,8,9,10,11,12,13,14,15,16,17,18,19,20} {
            \fill[cyan!50] (\h+0.05, 2.05) rectangle (\h+0.95, 2.95);
        }
        % VC row (y=1): hours 20-23
        \foreach \h in {20,21,22,23} {
            \fill[orange!50] (\h+0.05, 1.05) rectangle (\h+0.95, 1.95);
        }
        % Vacuum row (y=0): hours 5,16,19
        \foreach \h in {5,16,19} {
            \fill[purple!40] (\h+0.05, 0.05) rectangle (\h+0.95, 0.95);
        }

        % Draw grid - Optimized (offset by 26)
        \draw[very thin, gray] (26, -0.5) grid[step=1] (50, 5.5);

        % Y-axis labels (right side)
        \node[font=\tiny, anchor=east] at (25.8, 4.5) {WM};
        \node[font=\tiny, anchor=east] at (25.8, 3.5) {Heater};
        \node[font=\tiny, anchor=east] at (25.8, 2.5) {AC};
        \node[font=\tiny, anchor=east] at (25.8, 1.5) {VCharger};
        \node[font=\tiny, anchor=east] at (25.8, 0.5) {Vacuum};

        % OPTIMIZED SCHEDULE - fill ON cells
        % WM row: hours 0,1,2,3,4,6,7
        \foreach \h in {0,1,2,3,4,6,7} {
            \pgfmathsetmacro{\x}{\h+26}
            \fill[blue!50] (\x+0.05, 4.05) rectangle (\x+0.95, 4.95);
        }
        % Heater row: hours 22,23
        \foreach \h in {22,23} {
            \pgfmathsetmacro{\x}{\h+26}
            \fill[red!50] (\x+0.05, 3.05) rectangle (\x+0.95, 3.95);
        }
        % AC row: hours 0-17, 22, 23
        \foreach \h in {0,1,2,3,4,5,6,7,8,9,10,11,12,13,14,15,16,17,22,23} {
            \pgfmathsetmacro{\x}{\h+26}
            \fill[cyan!50] (\x+0.05, 2.05) rectangle (\x+0.95, 2.95);
        }
        % VC row: hours 0,2,3,4
        \foreach \h in {0,2,3,4} {
            \pgfmathsetmacro{\x}{\h+26}
            \fill[orange!50] (\x+0.05, 1.05) rectangle (\x+0.95, 1.95);
        }
        % Vacuum row: hours 1,22,23
        \foreach \h in {1,22,23} {
            \pgfmathsetmacro{\x}{\h+26}
            \fill[purple!40] (\x+0.05, 0.05) rectangle (\x+0.95, 0.95);
        }

        % Arrow between the two grids
        \draw[-{Stealth[length=4mm]}, ultra thick, green!50!black] (24.5, 2.5) -- node[above, font=\scriptsize\bfseries]{MILP} (25.5, 2.5);

        % Legend
        \node[font=\tiny] at (12, -2.8) {Background: \colorbox{green!15}{Off-Peak} \colorbox{yellow!15}{Day} \colorbox{red!15}{Peak} \quad Colored cells = ON};
    \end{tikzpicture}
    \caption{24-hour schedule heatmap: Original (LSTM) vs.\ Optimized (MILP). The optimizer shifts loads from peak/day to off-peak bands.}
    \label{fig:heatmap}
\end{figure*}

\subsection{Policy Validation}
All 82 policies pass (100\% pass rate, Table~\ref{tab:validation}). Key highlights: all 48 hourly capacity checks pass; Heater runs only at hours 22--23 within Ollama-preferred \{18--23\}; 95\% of AC ON-hours overlap with hot/humid conditions; total savings of 37.0\% exceed the 20\% research threshold.

\begin{table}[H]
    \centering
    \caption{Policy validation results.}
    \label{tab:validation}
    \small
    \begin{tabular}{@{}lcc@{}}
        \toprule
        \textbf{Category} & \textbf{Pass/Total} & \textbf{Rate} \\
        \midrule
        Capacity           & 48/48 & 100\% \\
        Format             & 10/10 & 100\% \\
        LLM Preference     & 10/10 & 100\% \\
        LSTM Accuracy      & 5/5   & 100\% \\
        Cost Optimality    & 6/6   & 100\% \\
        Savings Threshold  & 1/1   & 100\% \\
        Weather Comfort    & 2/2   & 100\% \\
        \midrule
        \textbf{Total}     & \textbf{82/82} & \textbf{100\%} \\
        \bottomrule
    \end{tabular}
\end{table}

\subsection{Savings Distribution}
Fig.~\ref{fig:savings_pie} shows the distribution of the total LKR~995.70 savings across appliances. The VehicleCharger and AC contribute the most absolute savings due to their high power ratings and successful peak-to-off-peak shifting.

% ---- SAVINGS DISTRIBUTION PIE CHART ----
\begin{figure}[H]
    \centering
    \begin{tikzpicture}
        \pie[
            radius=1.8,
            text=pin,
            pin distance=0.4cm,
            font=\tiny,
            color={blue!40, red!40, cyan!40, orange!40, purple!30},
            explode={0, 0.08, 0.08, 0.08, 0},
        ]{
            6.1/WM: 60.6,
            24.1/Heater: 240.0,
            26.4/AC: 262.8,
            32.3/VCharger: 321.2,
            11.2/Vacuum: 111.1
        }
    \end{tikzpicture}
    \caption{Savings distribution by appliance (\% of total LKR~995.70).}
    \label{fig:savings_pie}
\end{figure}

% ==============================================================================
% 5. DISCUSSION
% ==============================================================================
\section{Discussion}

\subsection{Interpretation of Results}
The 37.0\% cost reduction demonstrates significant potential for AI-driven demand-side management under LECO TOU tariffs. Savings are primarily from shifting high-power appliances (Heater: 64.5\%, VehicleCharger: 52.5\%) from peak/day to off-peak bands, exploiting the threefold price differential.

The integration of weather-based comfort ensures optimization does not compromise well-being. The AC, despite running 20/24~hours, achieves 95\% overlap with comfort-critical hours through the $-100$ bias term. The LLM successfully parsed \textit{``Heater need to turn only night times''} into structured constraints, demonstrating that natural-language user interaction is viable within an optimization pipeline.

The complete data flow---sensors $\rightarrow$ MQTT $\rightarrow$ LSTM $\rightarrow$ LangGraph Agent (MQTT TOU + Weather API + Firebase preferences + Ollama LLM) $\rightarrow$ MILP solver $\rightarrow$ Firebase $\rightarrow$ Flutter app---operates as a cohesive pipeline with the edge controller enabling real-time grid awareness.

\subsection{Comparison with Literature}
The 37\% savings compare favorably with Pipattanasomporn~et~al.~\cite{pipattanasomporn2012} (21\% peak reduction) and Garifi~et~al.~\cite{garifi2018} (15--25\% savings), neither of which incorporated demand forecasting, LLM preference parsing, or comfort-aware scheduling.

\subsection{Limitations}
\begin{enumerate}
    \item \textbf{Simulated sensors:} Current evaluation uses simulated data; real smart-plug deployment is needed.
    \item \textbf{Appliance set:} Five appliances; scaling to larger sets with renewables/storage requires extended MILP formulations.
    \item \textbf{LSTM architecture:} Fixed 24-step window; Transformer models may improve accuracy.
    \item \textbf{LLM robustness:} Tested with limited instruction set; diverse input evaluation needed.
    \item \textbf{Single household:} Multi-household coordination via federated edge controllers is future work.
\end{enumerate}

% ==============================================================================
% 6. CONCLUSION
% ==============================================================================
\section{Conclusion and Future Work}

This thesis presented an AI-based HEMS integrating LSTM demand forecasting (trained on 10,080 samples, 93.3\% scheduling accuracy), MQTT edge controllers for live LECO TOU tariffs, Ollama LLM for natural-language preference parsing, and MILP optimal scheduling---orchestrated by a LangGraph agent with automatic fallback. The system achieves \textbf{37.0\% daily cost reduction} (LKR~2,692.80~$\rightarrow$~1,697.10) with 82/82 validation policies passing (100\% rate) across capacity, cost, LSTM, LLM, and comfort categories.

These findings demonstrate that: (i)~LSTM models accurately predict appliance demand for scheduling; (ii)~MQTT edge controllers enable real-time grid-aware decisions; (iii)~locally hosted LLMs reliably parse natural-language preferences; (iv)~MILP produces provably optimal schedules respecting all constraints; and (v)~comprehensive policy validation provides deterministic correctness guarantees.

Future work will focus on physical deployment with smart plugs, integration of solar PV and battery storage, Transformer-based forecasting, multi-household coordination, and longitudinal user studies.

% ==============================================================================
% REFERENCES (Page 10)
% ==============================================================================
\onecolumn
\section*{References}
\addcontentsline{toc}{section}{References}

\begin{thebibliography}{13}
\small

\bibitem{leco2024}
Lanka Electricity Company (Pvt) Ltd, ``LECO Domestic Time-of-Use Tariff Schedule,'' 2024. [Online]. Available: \url{https://www.leco.lk}

\bibitem{siano2014}
P.~Siano, ``Demand response and smart grids---A survey,'' \textit{Renew. Sustain. Energy Rev.}, vol.~30, pp.~461--478, 2014.

\bibitem{zhou2016}
B.~Zhou \textit{et~al.}, ``Smart home energy management systems: Concept, configurations, and scheduling strategies,'' \textit{Renew. Sustain. Energy Rev.}, vol.~61, pp.~30--40, 2016.

\bibitem{beaudin2015}
M.~Beaudin and H.~Zareipour, ``Home energy management systems: A review of modelling and complexity,'' \textit{Renew. Sustain. Energy Rev.}, vol.~45, pp.~318--335, 2015.

\bibitem{pipattanasomporn2012}
M.~Pipattanasomporn, M.~Kuzlu, and S.~Rahman, ``An algorithm for intelligent home energy management and demand response analysis,'' \textit{IEEE Trans. Smart Grid}, vol.~3, no.~4, pp.~2166--2173, 2012.

\bibitem{hochreiter1997}
S.~Hochreiter and J.~Schmidhuber, ``Long short-term memory,'' \textit{Neural Comput.}, vol.~9, no.~8, pp.~1735--1780, 1997.

\bibitem{garifi2018}
K.~Garifi, K.~Baker, B.~Touri, and D.~Christensen, ``Stochastic model predictive control for demand response in a home energy management system,'' in \textit{Proc. IEEE PESGM}, 2018, pp.~1--5.

\bibitem{brown2020}
T.~Brown \textit{et~al.}, ``Language models are few-shot learners,'' in \textit{Adv. Neural Inf. Process. Syst.}, vol.~33, 2020, pp.~1877--1901.

\bibitem{naik2017}
N.~Naik, ``Choice of effective messaging protocols for IoT systems: MQTT, CoAP, AMQP and HTTP,'' in \textit{Proc. IEEE ISSE}, 2017, pp.~1--7.

\bibitem{kong2019}
W.~Kong \textit{et~al.}, ``Short-term residential load forecasting based on LSTM recurrent neural network,'' \textit{IEEE Trans. Smart Grid}, vol.~10, no.~1, pp.~841--851, 2019.

\bibitem{kim2019}
T.-Y.~Kim and S.-B.~Cho, ``Predicting residential energy consumption using CNN-LSTM neural networks,'' \textit{Energy}, vol.~182, pp.~72--81, 2019.

\bibitem{suganthi2012}
L.~Suganthi and A.~A.~Samuel, ``Energy models for demand forecasting---A review,'' \textit{Renew. Sustain. Energy Rev.}, vol.~16, no.~2, pp.~1223--1240, 2012.

\bibitem{king2023}
J.~King, ``LLM-powered smart home assistants: A survey,'' \textit{arXiv preprint}, 2023.

\end{thebibliography}

\end{document}
